% !PRINTED 660
\weberpage{660}

%<<CONT
\weq{3}{\begin{aligned}
\eta &= c + \varrho\eta_1, \\
\eta_1 &= c_1 + \varrho\eta_2, \\
&\;\cdots \\
\eta_{r-1} &= c_{r-1} + \varrho\eta_r,
\end{aligned}}

woraus sich ergibt:

\weq{4}{\eta = c + c_1\varrho + c_2\varrho^2 + \cdots + c_{r-1}\varrho^{r-1} + \eta_r\varrho^r.}

Die Reihe der Konstanten $c, c_1, c_2, \ldots, c_{r-1}$ ist durch
$\eta$ vollständig bestimmt und $\eta_1, \eta_2, \ldots, \eta_r$ ist
eine Reihe von Funktionen in $\Omega$, deren Nenner nicht durch $\pi$
teilbar ist.

Man könnte diesen Ausdruck die \gesperrt{Taylorsche Entwickelung} der Funktion
$\eta$ nach Potenzen von $\varrho$ nennen. Eine ähnliche Entwickelung
läßt sich auch für eine Funktion $\eta$ geben, deren Nenner durch
$\pi$ teilbar ist, nur daß dabei auch negative Potenzen auftreten:
Man kann nämlich eine Potenz $\varrho^s$ so bestimmen, daß der Nenner
von $\varrho^s\eta$ nicht mehr durch $\pi$ teilbar ist, und wenn man
auf dieses Produkt die Entwickelung (4) anwendet, so folgt:

\weq{5}{\eta = c\varrho^{-s} + c_1\varrho^{-s+1} + \cdots + c_{r-1}\varrho^{r-s-1} + \eta_r\varrho^{r-s}.}

\begin{center}
§ 186. \textbf{\gesperrt{Birationale Transformation.}}
\end{center}

Wir kommen jetzt zu einem Gegenstand, der das Gebiet der
algebraischen Funktionen vorzugsweise von dem der algebraischen
Zahlen unterscheidet, und der Grund dafür ist, daß die Theorie der
ganzen Funktionen und Funktionale für die Funktionen bei weitem
nicht den Charakter der Invarianz hat, wie bei den Zahlen; das ist
die birationale Transformation.

Ist $z_1$ irgend eine nicht konstante Funktion des Körpers $\Omega$
und $N(t - z_1)$ die $f$te Potenz einer irreduziblen Funktion $e$ten
Grades (§ 172, 6.), so besteht zwischen $z$ und $z_1$ eine Gleichung,
die in bezug auf $z_1$ vom $e$ten Grade ist, deren Grad in bezug auf
$z$, wenn Nenner und überflüssige Faktoren weggeschafft sind, mit
$e_1$ bezeichnet werde. Diese Gleichung sei

\weq{1}{G\bigl(\overset{e}{z_1},\, \overset{e_1}{z}\bigr) = 0.}

Unter den verschiedenen Gleichungen dieser Form gibt es nur eine
irreduzible und diese ist sowohl in bezug auf $z$ als in bezug auf
$z_1$ von möglichst niedrigem Grade (Bd. I, § 20).

Ist $\theta$ eine primitive Funktion des Körpers $\Omega$, so ist
nach § 172, (12)

\weq{2}{z_1^h\theta^k \qquad
\begin{aligned}
h &= 0, 1, \ldots, e-1 \\
k &= 0, 1, \ldots, f-1
\end{aligned}}
%CONT>>
